\begin{annexitem}{annex:p19.rat63-or-if}{thm:p19-rat63-or-if}{Truth equivalence, probability ordering and the certain-premise case}
\begin{proof}
In bivalent semantics, \((\neg A)\to B\) is false exactly
when \(\neg A\) is true and \(B\) false, namely the assignment
\((A,B)=(0,0)\). The inclusive disjunction is false exactly there
as well. Hence their truth sets coincide and each entails the other.
This proof fixes the material reading; it does not establish the
interpretation of an experimental sentence.

The disjunction event is the disjoint union of \(A\) and
\(\neg A\cap B\). With \(a=P(A)\) and \(P(\neg A)=1-a>0\),
the definition of conditional probability gives
\(P(\neg A\cap B)=(1-a)q\). Therefore
\(P(A\lor B)=a+(1-a)q\), and subtracting \(q\) gives
\(a(1-q)\ge0\), since \(a,q\in[0,1]\). Thus the disjunction's
degree cannot fall below the stipulated conditional assessment \(q\).
This proves the indicated one-premise p-validity direction, not equality
between a conditional-probability assessment and a Boolean event.

The four nonnegative masses in the counterexample sum to one.
\(P(\neg A)=.04+.16=.20>0\), so \(q=.04/.20=.20\).
The disjunction has mass \(.4+.4+.04=.84\). Hence the reverse
confidence-nondecrease condition fails for this admissible law.

If the disjunction is certain, then
\(1=a+(1-a)q\), or \((1-a)(1-q)=0\). Positive probability
of \(\neg A\) makes \(1-a>0\), so \(q=1\).
Without that positivity the ordinary ratio is undefined; no value
is assigned by this proof. The counterexample to nondecrease over
uncertain premises therefore does not refute the certain-premise result.
\end{proof}
\end{annexitem}

\begin{annexitem}{annex:p19.rat63-compatibility}{thm:p19-rat63-compatibility}{A common possibility is not a universal consequence}
\begin{proof}
The assignment \((A,B)=(1,1)\) satisfies both \(A\lor B\)
and \(A\to B\), so the intersection of their truth sets is
nonempty. The assignment \((1,0)\) satisfies the disjunction
and falsifies the material conditional. Thus it belongs to
\(\mathcal T(P_0)\setminus\mathcal T(Q_0)\), proving the failure
of inclusion. Existential compatibility is insufficient for the
universal condition defining entailment. This calculation uses fully
specified valuations, not the incomplete initial cognitive models.
\end{proof}
\end{annexitem}

\begin{annexitem}{annex:p19.rat63-observe-intervene}{thm:p19-rat63-observe-intervene}{Evidence at one world and action at another}
\begin{proof}
In the factual model both \(B\) and \(C\) equal \(U\).
The evidence \(B=1\) has probability \(1/2\) and restricts
the posterior background to \(U=1\). In that factual world
\(C=1\), so the conditional probability of \(C=0\) is zero.

Replacing \(B\)'s equation by the constant zero while preserving
the same posterior unit leaves \(C=B\) as its equation and hence
gives \(C[B=0]=0\). This holds with posterior probability one.
The replacement is evaluated at a different world label; it does
not change what factual evidence was supplied. In contrast no factual
outcome satisfies \(B=1\land B=0\); its probability is zero.
The ordinary conditional ratio on that event is therefore undefined.
All calculations concern the supplied equations, not a measured
representation or an estimate from a psychological sample.
\end{proof}
\end{annexitem}

\begin{annexitem}{annex:p19.rat63-alternates}{thm:p19-rat63-alternates}{Exhaustiveness, category complements and underdetermination}
\begin{proof}
Let \(S\) be the domain of the single category and let \(R,P,L\)
be the events for roses, poppies and lilies. In the binary context,
\(R\cup P=S\) and \(R\cap P=\varnothing\). Every point
outside \(R\) lies in \(P\), and every point in \(P\) lies
outside \(R\), so \(S\setminus R=P\).

In the three-category context, disjointness and exhaustiveness similarly
give \(S\setminus R=P\cup L\). The stipulated possible poppy
and lily categories supply two distinct states consistent with not-roses;
one falsifies lilies, the other falsifies poppies. Thus the negation
does not determine either individual alternate. If a named category
is impossible in a different context, that is a further restriction
and not the three-possible-category case just stated.

Without exhaustiveness, adjoin an unnamed category \(Z\) disjoint
from all the named ones. A state in \(Z\) satisfies not-roses but
not poppies-or-lilies. This is the required counterexample to the
unqualified implication. The category equivalences say nothing about
whether an agent actually represented alternates or symbolic negation.
They also add no converse or biconditional assumption to a material
\(A\to B\); the existing AC/DA countermodels therefore remain.
\end{proof}
\end{annexitem}
